\documentclass[14pt]{mmcs_article}

\usepackage[utf8]{inputenc}
\usepackage[T2A]{fontenc}
\usepackage[english,russian]{babel}
\usepackage[skip=10pt, indent=1.25cm]{parskip}

\usepackage{algorithm}
\usepackage{algpseudocode}
\usepackage{tabularray}
\usepackage{amsthm}
\UseTblrLibrary{amsmath}
\usepackage{amssymb}
\usepackage{multicol}
\usepackage{setspace}

\usepackage{csquotes}
\usepackage{hyperref}

% https://tex.stackexchange.com/questions/37797/theorem-environment-line-break-after-label
\newtheoremstyle{break}
  {}          % Space above, empty = `usual value'
  {}          % Space below
  {\mdseries} % Body font
  {}          % Indent amount (empty = no indent, \parindent = para indent)
  {\bfseries} % Thm head font
  {.}         % Punctuation after thm head
  { }
  {}          % Thm head spec
\theoremstyle{break}
\newtheorem{definition}{Определение}[section]
\newtheorem{theorem}{Теорема}[section]
\newtheorem{lemma}{Лемма}[theorem]
\newtheorem{remark}{Замечание}
  %{\newline}  % Space after thm head: \newline = linebreak

\usepackage[backend=biber]{biblatex}
\addbibresource{main.bib}

 % For aligned numeric columns in comparison table
\usepackage{siunitx}
\sisetup{
  detect-mode,
  tight-spacing = true,
  table-number-alignment = center,
  round-mode = places,
  round-precision = 2,
  per-mode = symbol
}

\hypersetup{
    colorlinks=true,
    filecolor=magenta,
    linkcolor=blue,
    citecolor=cyan,
    urlcolor=cyan,
}
\urlstyle{same}

\begin{document}

\renewcommand{\contentsname}{Оглавление}
\tableofcontents
\newpage

\section{Введение}

% https://csrc.nist.gov/News/2022/pqc-candidates-to-be-standardized-and-round-4
% Не про KEM
В 2024 году завершилось соревнование \textquote{Post-Quantum Cryptography Standartization}, проводимое NIST \cite{nist-post-quantum-crypto-standardization}. Целью соревнования было включить в список стандартов NIST новые криптосистемы, устойчивые к атакам со стороны квантовых компьютеров. В частности, в отборе учавстовали механизмы инкапсуляции ключей.

Механизм инкапсуляции ключей (KEM, Key Encapsulation Mechanism)~--- это набор алгоритмов, которые могут использоваться двумя сторонами для безопасного установления общего секретного ключа через публичный канал \cite{nist-recommendations-for-kems}. Общий секретный ключ затем может быть использован с традиционными симметричными криптографическими алгоритмами для, например, шифрования или аутентификации.

В четвертый раунд соревнования прошли четыре KEM: BIKE, Classic McEliece, HQC, SIKE \cite{nist-round-4-report}. В данной работе рассматривается BIKE --- криптосистема, основанная на QC-MDPC кодах \cite{bike-website}. По итогам четвёртого раунда NIST отдал предпочтение HQC как резервному стандарту, отметив, что анализ частоты ошибок декодирования (DFR) для BIKE пока менее зрелый \cite{nist-pqc-status-report-on-fourth-round}. Тем не менее BIKE остаётся активно исследуемой криптосистемой: она обеспечивает наименьший размер ключей среди всех кандидатов и допускает эффективную программную реализацию. 

\paragraph{Структура работы.} Раздел \ref{sec2} вводит определения и обозначения, используемые в дальнейшем. В разделе \ref{sec3} будут показаны общие подходы к вычислению циклических сверток, а в разделе \ref{sec4}~-- общие подходы к обращению многочленов. В разделе \ref{sec5} описана оптимизированная реализация циклических сверток над $F_4$. В разделе \ref{sec6} будет показано, как перенести оптимизированный алгоритм обращения многочлена для $F_2$ на поле $F_4$. Раздел \ref{sec7} содержит подробности о программной реализации и результаты вычислительных экспериментов. Выводы из проведенной работы содержатся в разделе \ref{sec8}.

\section{Предварительные сведения}
\label{sec2}

\subsection{Определения и обозначения}

Через $F_p$ обозначим конечное поле мощности $p$, где $p$ -- простое число. Через $F_p[x]$ обозначим кольцо многочленов с коэффициентами из этого поля. Пусть $f(x) \in \mathbb{F}_p[x]$ --- многочлен степени $n$. Через $R = \frac{\mathbb{F}_p[x]}{f(x)}$ обозначим \emph{факторкольцо многочленов} по идеалу, порожденному $f$. Оно состоит из многочленов степени меньше $n$ с коэффициентами из $\mathbb{F}_p$, а умножение выполняется по модулю $f(x)$. В криптосистемах на QC-MDPC кодах используется $f(x) = x^n - 1$, так что $x^n \equiv 1$. Из-за этого умножение многочленов в этих криптосистемах может быть реализовано \emph{циклической сверткой}.

Циклическая свёртка двух последовательностей $x[n]$ и $h[n]$ длины $N$ определяется как

\begin{equation}
    (x \circledast h)[n] = \sum_{m=0}^{N-1} x[m]\, h[(n - m) \bmod N],
    \label{eq:cycconv}
\end{equation}

Линейным $[n, k]_q$ кодом называют $k$-мерное подпространство $\mathbb{F}_q^n$. Элементы кода называются кодовыми словами. Линейный код можно задать проверочной матрицей $H \in \mathbb{F}_q^{(n-k)\times n}$: элемент $c \in \mathbb{F}^n$ является кодовым словом тогда и только тогда, когда $H c^T = 0$.

Поле $\mathbb{F}_4 = \{0, 1, \alpha, \alpha+1\}$ является расширением $\mathbb{F}_2$ степени~2, задаваемым неприводимым многочленом $x^2 + x + 1$. Для хранения элементов из $\mathbb{F}_4$ нужны всего 2 бита, что позволяет выполнять операции над большим количеством элементов одновременно с использованием побитовых операций.

\paragraph{Малая теорема Ферма.}
В конечном поле $\mathbb{F}_{q^m}$ для любого ненулевого элемента $a$ выполняется $a^{q^m - 1} = 1$.

\paragraph{Китайская теорема об остатках.}
Если $f(x) = f_1(x) \cdots f_l(x)$, где многочлены $f_i$ попарно взаимно просты, то

\begin{equation}
    \frac{\mathbb{F}_q[x]}{f(x)} \cong \frac{\mathbb{F}_q[x]}{f_1} \oplus \frac{\mathbb{F}_q[x]}{f_2} \oplus \cdots \oplus \frac{\mathbb{F}_q[x]}{f_l}
\end{equation}

\subsection{Быстрое возведение в степень}

Для возведения многочлена в степень будет использоваться алгоритм быстрого возведения в степень.

\begin{algorithm}[!h]
\caption{Быстрое возведение в степень}
\begin{algorithmic}[1]
\State $e \leftarrow S$
\State $\mathit{result} \leftarrow 1$
\State $\mathit{base} \leftarrow a$
\While{$e > 0$}
    \If{$e$ нечётно}
        \State $\mathit{result} \leftarrow \mathit{result} \cdot \mathit{base} \bmod (x^n - 1)$
    \EndIf
    \State $\mathit{base} \leftarrow \mathit{base}^2 \bmod (x^n - 1)$
    \State $e \leftarrow \lfloor e / 2 \rfloor$
\EndWhile
\State \textbf{return} $\mathit{result}$
\end{algorithmic}
\end{algorithm}

\subsection{Дискретное преобразование Фурье}
\label{dpf}

Пусть $x[n]$, $n = 0, 1, \ldots, N-1$, — конечная последовательность
комплексных чисел. \emph{Дискретным преобразованием Фурье} (ДПФ) называется
отображение

\begin{equation}
    X[k] = \sum_{n=0}^{N-1} x[n]\, W_N^{kn}, \qquad k = 0, 1, \ldots, N-1,
    \label{eq:dft}
\end{equation}

\noindent где $W_N = e^{-j2\pi/N}$ — примитивный корень степени $N$ из
единицы. Обратное ДПФ (ОДПФ) определяется формулой

\begin{equation}
    x[n] = \frac{1}{N} \sum_{k=0}^{N-1} X[k]\, W_N^{-kn}, \qquad
    n = 0, 1, \ldots, N-1.
    \label{eq:idft}
\end{equation}

Прямое вычисление требует $\mathcal{O}(N^2)$ арифметических операций. Алгоритмы быстрого преобразования Фурье сокращают эту сложность до $\mathcal{O}(N \log N)$, что делает ДПФ практически применимым при больших $N$.

Важным практическим свойством ДПФ возможность вычислять с его помощью циклические свертки. Результат циклической свертки можно вычислить как поточечное произведение спектров:

\begin{equation}
    \mathcal{F}\{x \circledast h\}[k] = X[k] \cdot H[k],
\end{equation}

\noindent где $X[k]$ и $H[k]$ \-- ДПФ последовательностей $x[n]$ и $h[n]$
соответственно. Таким образом, циклическая свёртка вычисляется за три шага:

\begin{equation}
    x \circledast h = \mathcal{F}^{-1}\bigl(\mathcal{F}(x) \cdot \mathcal{F}(h)\bigr),
\end{equation}

\noindent что при использовании быстрого преобразования Фурье требует $\mathcal{O}(N \log N)$ операций вместо $\mathcal{O}(N^2)$ при прямом вычислении.

\subsection{Криптосистемы, основанные на QC-MDPC кодах}

MDPC (Medium density parity check) код --- это линейный $[n, k]_q$ код, у которого проверочная матрица $H$ имеет вес строк $\omega \approx \sqrt{n}$.

Говорят, что квадратная матрица $A$ является циркулянтом, если в ней каждая строка получается циклическим сдвигом предыдущей строки вправо:

\begin{equation}
A = \begin{pmatrix}
    1 & 1 & 0 \\
    0 & 1 & 1 \\
    1 & 0 & 1
\end{pmatrix}
\end{equation}

Из-за такой структуры вместо хранения всей матрицы можно хранить только её первую строку. Матрицам циркулянтам $A$ $B$ можно сопоставить многочлен... Если представить эту строку как многочлен, то результат умножения матриц-циркулянтов будет таким же, как умножение в факторкольце соответствующих многочленов.

MDPC код, обладающий блочно-циркулянтной проверочной матрицей, называется квази-циклическим (QC-MDPC, Quasi-Cyclic Medium Density Parity Check).

\begin{equation}
H = \begin{+pmatrix}[colspec={rrr|rrr}]
    1 & 0 & 1 & 0 & 0 & 1 \\
    1 & 1 & 0 & 1 & 0 & 0 \\
    0 & 1 & 1 & 0 & 1 & 0
\end{+pmatrix}
\end{equation}

Рассмотрим, как в криптосистеме на QC-MDPC кодах реализуются эти операции.

Секретный ключ $H$ --- проверочная матрица случайного QC-MDPC $[n = ln', (l-1)n']_q$-кода, представленная как

\begin{equation}
    H = \begin{pmatrix} H_1 \mid H_2 \mid \ldots \mid H_l \end{pmatrix},
\end{equation}

где каждый циркулянтный блок $H_i$ размера $n' \times n'$ имеет вес строки $\gamma$ (равный примерно $\sqrt{n}$). Параметр $n'$ --- простое число. В криптосистеме, разработанной в этой работе параметр $l = 2$, и дальше всегда будет предполагаться что $l = 2$. Публичный ключ $\widetilde{H}$ --- систематическая форма $H$:

\begin{equation}
    \widetilde{H} = H_1^{-1} H
    = \begin{pmatrix} I_{n'} \mid H_1^{-1} H_2 \mid \ldots \mid H_1^{-1} H_l \end{pmatrix}.
\end{equation}

Поскольку произведение матриц-циркулянтов есть матрица-циркулянт, $\widetilde{H}$ также является блочно-циркулянтной. Публичный ключ представляется первыми строками блоков $H_1^{-1} H_i$, $i \in \{2, \ldots, l\}$. Таким образом, при $l = 2$, публичный ключ является многочленом степени $n' - 1$.

Общий секрет --- вектор ошибок $e \in \mathbb{F}_q^n$ веса $t$, а шифртекст --- его синдром $\widetilde{s} = \widetilde{H} e^T$. Для расшифрования вычисляется частный синдром $s = H e^T = H_1 \widetilde{s}^T$, который подаётся на вход MDPC-декодеру.

В BIKE параметр $l$ выбран равным 2, и все операции реализованы над $\mathbb{F}_2$. Оптимизированная реализация за константное время представлена в \cite{bike-reference-implementation}. В данной работе разрабатываются быстрые алгоритмы, поддерживающие поле $\mathbb{F}_4$, а также произвольные поля $\mathbb{F}_p$ для простого $p$.

\paragraph{Шифрование с помощью KEM} KEM реализует три алгоритма: \texttt{KeyGen}, \texttt{Encaps} и \texttt{Decaps}. Алиса генерирует публичный и приватный ключ с помощью \texttt{KeyGen} и отправляет публичный ключ Бобу. Боб генерирует общий секретный ключ и шифрует его публичным ключом через \texttt{Encaps}, затем отправляет полученный шифртекст Алисе. Алиса расшифровывает его приватным ключом через \texttt{Decaps} и получает тот же общий секрет. Таким образом, участники договорились об общем секрете, не передавая его явно по открытому каналу.

\subsection{QC-MDPC Декодер}

Для декодирования шифртекста в криптосистемах MDPC используется итеративный декодер, восстанавливающий вектор ошибок $e$ по синдрому $s$ и матрице $H$. Для QC-MDPC кодов применяются декодеры класса \emph{symbol flipping}, которые являются расширением bit-flipping декодера.

\paragraph{Bit-flipping decoder}

На каждой итерации декодера для каждой позиции $j \in [0, n - 1]$ вычисляется счётчик

\begin{equation}
    \sigma_j = \bigl|\{ i : H_{ij} = 1,\; s_i = 1 \}\bigr|
\end{equation}

Вычисленное $\sigma_j$ --- количество проверочных уравнений, которым позиция $j$ не удовлетворяет. Если $\sigma_j$ превышает пороговое значение $threshold$, бит $e_j$ переключается и синдром пересчитывается. Итерации повторяются до обнуления синдрома или до достижения максимального числа итераций.

\section{Циклических свертки в $\frac{F_p[x]}{x^n - 1}$}
\label{sec3}

Самой частой операцией в криптосистеме является циклическая свертка. Напомним, что матрицы циркулянты, которые являются основными объектами в QC-MDPC криптосистеме, могут быть представлены своей первой строкой, как многочлен из факторкольца $\frac{\mathbb{F}_q[x]}{x^n - 1}$. В свою очередь умножение многочленов из этого факторкольца эквивалентно циклической свертке из-за $x^n \equiv 1$. Так как эта операция выполняется часто, имеет смысл применить к ней несколько базовых оптимизаций.

\begin{algorithm}[!htpb]
\caption{Циклическая свёртка $c = a \circledast b$}
\begin{algorithmic}[1]
    \For{каждой позиции $i$}
        \For{каждой позиции $j$}
            \State $c_i \leftarrow c_i + a_j b_{(i - j) mod\ n}$
        \EndFor
    \EndFor
\end{algorithmic}
\end{algorithm}

Операция взятия остатка от деления вычислительно дороже чем обычные арифметические операции. От неё можно просто избавиться, расщепив цикл на два:

\begin{algorithm}[!htpb]
\caption{Циклическая свёртка $c = a \circledast b$ без взятия остатка от деления}
\begin{algorithmic}[1]
    \For{каждой позиции $i$}
        \For{каждой позиции $j \leq i$}
            \State $c_i \leftarrow c_i + a_j b_{i - j}$
        \EndFor
        \For{$j$ от $i + 1$ до $n - 1$}
            \State $c_i \leftarrow c_i + a_j b_{n + (i - j)}$
        \EndFor
    \EndFor
\end{algorithmic}
\end{algorithm}

\subsection{Разреженные свертки}

Напомним, что секретный ключ криптосистемы является матрицей с весом строки $\gamma$, который примерно равен $\sqrt{n}$. Это означает, что многочлен, описывающий эту матрицу, разреженный. Используя этот факт, можно выполнять операции только для ненулевых позиций:

\begin{algorithm}[!h]
\caption{Циклическая свёртка $c = a \circledast b$, где $a$ -- разреженный многочлен}
\begin{algorithmic}[1]
    \For{каждой позиции $i$}
        \For{каждой позиции $j$ такой, что $a_j \neq 0$}
            \State $c_i \leftarrow c_i + a_j b_{(i - j) mod\ n}$
        \EndFor
    \EndFor
\end{algorithmic}
\end{algorithm}

\subsection{Быстрое преобразование Фурье}

Циклическую свертку можно выполнять при помощи дискретного преобразования Фурье, как было отмечено в разделе \ref{dpf}. Однако это неприменимо в криптографическом контексте, из-за того что дискретное преобразование Фурье выполняется в поле комплексных чисел, и операции с числами с плавающей точкой не поддаются анализу так же хорошо, как целые числа.

\section{Обращения многочленов в $\frac{F_p[x]}{x^n - 1}$}
\label{sec4}

Генерация ключей требует нахождения обратного элемента $p^{-1} \in \frac{\mathbb{F}_q[x]}{x^n - 1}$, то есть требуется найти такой многочлен $v$, что:

\begin{equation}
    p \cdot v \equiv 1 \pmod{x^n - 1}.
\end{equation}

Это самая вычислительно дорогая операция. Основных алгоритмов два: расширенный алгоритм Евклида и подход, использующий малую теорему Ферма.

\subsection{Расширенный алгоритм Евклида в $\mathbb{F}_q[x]$}

Классический расширенный алгоритм Евклида для многочленов над полем $\mathbb{F}_q$ находит НОД двух многочленов $a, b \in \mathbb{F}_q[x]$ вместе с коэффициентами Безу $s, t$ такими, что

\begin{equation}
    s \cdot a + t \cdot b = \gcd(a, b).
\end{equation}

Применяя расширенный алгоритм Евклида к $a = x^n - 1$ и $b = p(x)$, получим

\begin{equation}
    s \cdot (x^n - 1) + t \cdot p(x) = \gcd(x^n - 1,\, p(x)).
\end{equation}

Теперь возьмем полученное тождество по модулю $(x^n - 1)$:

\begin{equation}
    t \cdot p(x) \equiv \gcd(x^n - 1,\, p(x)) \pmod{x^n - 1}.
\end{equation}

Таким образом, обратный элемент существует тогда и только тогда, когда $\gcd(x^n - 1, p(x)) = c,\quad c \in \mathbb{F}_q$, и в этом случае

\begin{align}
    t \cdot p(x) &\equiv c \pmod{x^n - 1} \\
    p^{-1}(x) &= c^{-1} t(x) \pmod{x^n - 1}
\end{align}

Ниже приведена реализация алгоритма в псевдокоде:

\begin{algorithm}
\caption{Расширенный алгоритм Евклида для нахождения $p^{-1}$, $p \in \mathbb{F}_q[x]/(x^n-1)$}
\begin{algorithmic}[1]
\State $a \leftarrow x^n - 1$
\State $b \leftarrow p$
\State $v_{\mathrm{old}} \leftarrow 0$
\State $v \leftarrow 1$
\Loop
    \State Вычислить частное $q \leftarrow \lfloor a / b \rfloor$ и остаток $a \leftarrow a - q \cdot b$
    \If{$a = 0$}
        \State \textbf{return} $b^{-1} \cdot v \bmod (x^n - 1)$
        \Comment{$b$ — константа, $v$ — искомый обратный}
    \EndIf
    \State $v_{\mathrm{new}} \leftarrow v_{\mathrm{old}} - q \cdot v \bmod (x^n - 1)$
    \State $(a,\, b) \leftarrow (b,\, a)$
    \State $(v_{\mathrm{old}},\, v) \leftarrow (v,\, v_{\mathrm{new}})$
\EndLoop
\end{algorithmic}
\end{algorithm}

\noindent Несколько деталей реализации:
\begin{itemize}
    \item Деление $\lfloor a / b \rfloor$ на шаге~6 выполняется стандартным делением «уголком» в $\mathbb{F}_q[x]$: на каждой итерации внутреннего цикла старший коэффициент $a$ обнуляется вычитанием подходящего кратного $b$, используя таблицу деления в $\mathbb{F}_q$.
    \item Произведение $q \cdot v \bmod (x^n-1)$ на шаге~10 реализовано циклической сверткой.
\end{itemize}

\subsection{Нахождение обратного элемента через малую теорему Ферма}

Пусть $R = \frac{\mathbb{F}_q[x]}{(x^n - 1)}$. Разложим $(x^n - 1)$ в произведение неприводимых многочленов $f_1, f_2, \ldots, f_l$ по китайской теореме остатков:

\begin{equation}
    R = \frac{\mathbb{F}_q[x]}{f_1} \oplus \frac{\mathbb{F}_q[x]}{f_2} \oplus \ldots \oplus \frac{\mathbb{F}_q[x}{f_l}
\end{equation}

Где каждое из $\frac{\mathbb{F}_q[x]}{f_i}$ изоморфно полю $F_{q^{deg(f_i)}}$. Тогда в этих полях выполняется малая теорема Ферма

\begin{equation}
    a^{q^{deg(f_i)} - 1} = 1
\end{equation}

И, эквивалентно:

\begin{equation}
    a^{q^{deg(f_i)} - 2} = a^{-1}
\end{equation}

Тогда в кольце $R$ по китайской теореме об остатках будет справедливо

\begin{equation}
    a^{lcm(q^{deg(f_1)} - 2, q^{deg(f_2)} - 2, \ldots, q^{deg(f_l)} - 2)} = a^{-1}
\end{equation}

Обозначим за $S$ степень $lcm(q^{deg(f_1)} - 2, q^{deg(f_2)} - 2, \ldots, q^{deg(f_l)} - 2)$.

Поскольку показатель $S$ экспоненциально велик, наивное возведение в степень неприемлемо. Применяется метод \emph{быстрого возведения в степень}: $S$ записывается в двоичном виде $S = \sum_{i=0}^{k} S_i 2^i$, и результат вычисляется за $O(\log S)$ умножений в $R$.

\subsection{Сравнение с расширенным алгоритмом Евклида}

Метод на основе малой теоремы Ферма требует $O(\log S)$ умножений многочленов (циклических сверток) в $R$, каждое из которых занимает $O(n^2)$ операций в $\mathbb{F}_q$ (или $O(n \log n)$ с использованием быстрого преобразования Фурье). Экспериментально было показано, что расширенный алгоритм Евклида выполняет порядка $O(n)$ итераций с $O(n)$ операциями каждая, то есть $O(n^2)$ суммарно. Таким образом, расширенный алгоритм Евклида предпочтительнее на практике.

\section{Оптимизация циклических сверток в $\frac{F_4[x]}{x^n - 1}$}
\label{sec5}

В разделе \ref{sec3} были показаны некоторые общие оптимизации циклических сверток, однако действительно значимый прирост можно получить только используя все возможности современных процессоров. Для этого можно представить свертку в $F_4$, как несколько сверток в $F_2$, для которых уже существуют оптимизированные реализации.

Представим $a \in F_4$ как два элемента из $F_2$, $a_r$  и $a_i$, $a = a_r + z a_i$. Тогда при умножении двух элементов получим:

\begin{equation}
    a b = (a_r b_r + z^2 a_i b_i) + z (a_r b_i + a_i b_r) \\
        = (a_r b_r + a_i b_i) + z (a_r b_i + a_i b_r + a_i b_i)
\end{equation}

Так, можно заметить что умножение в $F_4$ можно выполнить за счёт четырех умножений в $F_2$. Число умножений можно сократить до трех, воспользовавшись алгоритмом Карацубы.

Существует много оптимизированных алгоритмов, осуществляющих свертки в $F_2$. В данной работе была использована реализация из BIKE.

\section{Оптимизация обращений многочленов из $\frac{F_4[x]}{x^n - 1}$}
\label{sec6}

В разделе \ref{sec4} был сделан вывод о предпочтительности расширенного алгоритма Евклида для обращения многочлена. Однако у $F_4$ есть свойства, которые дают возможность реализовать значительно более быстрый алгоритм. Так, операция возведения в квадрат в $\mathbb{F}_4[x]/(x^n - 1)$ сводится к перестановке коэффициентов (циклическому сдвигу), аналогично случаю $\mathbb{F}_2$, поскольку $(a + b\alpha)^2 = a^2 + b^2 \alpha^2 = a^2 + b^2(\alpha + 1)$, что вычисляется дешевле, чем полное умножение. Это позволяет заменить значительную часть умножений в алгоритме быстрого возведения в степень операциями $k$-возведения в квадрат, следуя схеме алгоритма, использованного авторами BIKE~\cite{drucker-fast-polynomial-inversion}.

\paragraph{Утверждение о показателе $S$ для $\mathbb{F}_4$.}

Важным условием для применения алгоритма~\cite{drucker-fast-polynomial-inversion} является показатель степени $S$, в который возводится многочлен. К счастью, для простого $n$ показатели для $F_4$ и $F_2$ полностью совпадают, что позволяет перенести алгоритм без изменения количества умножений.

Покажем, почему показатели совпадают: пусть $n$ --- простое число, выбранное так, что $x^n - 1$ раскладывается в произведение трёх неприводимых многочленов над $\mathbb{F}_4$: $x^n - 1 = f_1 f_2 f_3$ с $\deg f_i = d_i$. Тогда показатель степени $S$ в кольце $\mathbb{F}_4[x]/(x^n - 1)$ совпадает с показателем, возникающим в случае $\mathbb{F}_2$:

\begin{equation}
    S = \mathrm{lcm}\!\bigl(4^{d_1} - 2,\; 4^{d_2} - 2,\; 4^{d_3} - 2\bigr) = 2^{n-1} - 2.
\end{equation}

Таким образом, при корректном выборе $n$, алгоритм~\cite{drucker-fast-polynomial-inversion} из BIKE переносится на случай $F_4$, с некоторыми изменениями, которые не оказывают существенного влияния на скорость выполнения.

\section{Результаты численных экспериментов}
\label{sec7}

В криптосистеме BIKE для достижения 192 бит классической стойкости или же 96 бит пост-квантовой безопасности (с учетом алгоритма Гровера) параметр $n$ был выбран равным 12323. Самыми дорогими операциями являются взятие обратного элемента в факторкольце $\frac{\mathbb{F}_q[x]}{x^n - 1}$ (нужно в \texttt{KeyGen} для публичного ключа $H^{-1}_1 H$), а также циклические свертки (произведение циркулянтов). Реализация выполнена на языке C++ и доступна в \cite{github-repo}.

\texttt{KeyGen}, \texttt{Encaps} и \texttt{Decaps}
На языке C++ были реализованы алгоритмы \texttt{KeyGen}, \texttt{Encaps} и \texttt{Decaps} в двух вариантах:

\begin{itemize}
    \item Наивная реализация для полей с порядком, равным произвольному простому числу.
    \item Ускоренная реализация для поля $F_4$.
\end{itemize}

Параметры криптосистемы: $n = 12323$, $\gamma = 71$, $t = 132$. Количество итераций декодера: 1. Параметры криптосистемы BIKE Level 3: $n = 24659$, $\gamma = 103$, $t = 199$. Количество итераций декодера: 4. \emph{Важно отметить, что поскольку в реализованной криптосистеме используется $F_4$, а в BIKE используется $F_2$. Поэтому для корректного сравнения в таблице параметр $n = 12323$ для криптосистемы над $F_4$ был умножен на 2 (24646)}.

Обе программы были собраны с максимальным уровнем оптимизации и с распараллеливанием с использованием инструкций AVX2. В таблице~\ref{tab:perf_gfx} приведено сравнение базовых операций между реализацией от авторов BIKE и реализацией в данной работе (GF4X). Операции gf2x\_mod\_inv и gf4x\_mod\_inv -- обращение многочлена, а gf2x\_mod\_mul и gf4x\_mod\_mul -- умножение многочленов (циклическая свертка).

В таблице~\ref{tab:perf} приведено сравнение производительноси трёх основных операций KEM: \texttt{KeyGen}, \texttt{Encaps} и \texttt{Decaps} между реализацией BIKE, реализацией над GF4, а также для наивной реализации.

\begin{table}[ht]
\centering
    \caption{Сравнение производительности между операциями над $F_2$ из BIKE Level 3, и реализацией этих же операций для $F_4$}
\label{tab:perf_gfx}
\begin{tabular}{lcccc}
\hline
Алгоритм & N & Операция & Время выполнения (мс) \\
\hline
BIKE Level 3 & 24659 & gf2x\_mod\_inv & 12.070 \\
             &       & gf2x\_mod\_mul & 0.571  \\
\hline
GF4X         & 24646 & gf4x\_mod\_inv & 9.551 \\
             &       & gf4x\_mod\_mul & 0.536  \\
\hline
\end{tabular}
\end{table}

\begin{table}[ht]
\centering
    \caption{Сравнение производительности между BIKE Level 3, реализацией для $F_4$ и наивной реализацией (GFPX, с $P = 5$)}
\label{tab:perf}
\begin{tabular}{lcccc}
\hline
Алгоритм & N & Операция & Время выполнения (мс) \\
\hline
BIKE Level 3 & 24659 & Keygen & 12.717 \\
             &       & Encaps & 0.730  \\
             &       & Decaps & 12.419 \\
\hline
GF4X         & 24646 & Keygen  & 12.587 \\
             &       & Encaps  & 0.721  \\
             &       & Decaps  & 4.607  \\
\hline
GFPX         & 24646 & Keygen  & 2009   \\
             &       & Encaps  & 829    \\
             &       & Decaps  & 22     \\
\hline
\end{tabular}
\end{table}

Важно отметить, что декодер, используемый в BIKE Level 3 выполняет 4 итерации, в то время как декодер из GF4X выполняет 1 итерацию. Поэтому для корректного сравнения нужно рассматривать значение $12.419 / 4 = 3.10475$ для Decaps из BIKE Level 3.

\section{Выводы}
\label{sec8}

Полученная реализация криптосистемы для полей $F_4$ сравнима в производительности по сравнению с реализацией от авторов BIKE для $F_2$. Эта реализация может иметь более хороший параметр DFR \cite{vedenev-paper}, а также может быть более сложной для взлома, ведь быстрых алгоритмов для $F_2$ больше чем для $F_4$.

\printbibliography

\end{document}
